% 前言

\section*{前言}
\addcontentsline{toc}{section}{前言}

从事计算机教学这些年，常有学生向编者提出类似的困惑：云计算的技术层出不穷，究竟要学到什么程度才算真正入门？
他们熟悉一个又一个新名词——容器、服务网格、无服务器计算，却总觉得跟不上技术迭代的节奏，越学越焦虑。

这样的困惑并非没有来由。
云计算的技术栈层次多、更新快，如果只是顺着名词逐一去记，很容易陷入"学了又忘、忘了又学"的循环，也很难判断一个新出现的技术究竟解决了什么问题、值不值得投入精力。
问题的关键不在于掌握了多少个名词，而在于有没有一条由真实问题驱动的线索，能够把零散的知识串成一条因果链。
理解了系统在规模、并发、状态、故障和资源的约束下如何演化出相应的机制，再去看任何一个技术名词，都能很快定位它所处的层次、所要解决的问题，以及为此付出的代价。

本书正是沿着这条线索组织内容的。
它不从某个工具或协议开始讲解，而是从一个最小的在线游戏原型出发，让系统在规模压力的推动下逐步演进。网络连通之后，并发问题随之而来；单机容量接近上限，就需要拆分到多个节点；节点数量增长，故障与一致性又成为新的挑战；最后将整套系统部署到云平台上，在弹性与成本之间反复权衡。每一个机制都对应前一阶段暴露的问题，每一种选择都有其适用条件和代价。

学完本书，读者获得的将不是又一张名词清单，而是一套分析和构建在线系统的方法。
从游戏案例中形成的判断，也可以迁移到在线协作、电商平台、大模型服务等其他系统。
这是本书编写的出发点，也是编者最希望读者带走的东西。

\section*{为什么写这本书}

云计算的书大体有两种。
一种把它讲成概念体系，IaaS、PaaS、SaaS 的定义滴水不漏，学生背完了 CAP 的三个字母，却写不出一个扛得住上千并发连接的服务器。
另一种把它讲成操作手册，照着文档点控制台、贴配置文件，离开具体工具便寸步难行。
前者有理无物，后者有物无理。
而云计算偏偏是一门理与物须臾不可分的学问：不懂原理，配置就是黑魔法；不动手实践，原理就是口头禅。
本书希望补足的，正是从问题、机制到实现和验证的完整链路。

这条链路之所以难建立，是因为云计算的难是一种组合的难。
拆开看，网络、操作系统、分布式系统这些零件，学生在先修课里都学过，也都可以单独找书来学。
但学过 TCP 三次握手，不等于知道一万并发连接下 listen 队列会溢出、半开连接该如何探测；学过进程线程，不等于知道几百个协程抢一把锁时吞吐为什么会雪崩；学过一致性定理，不等于知道一次跨服交易究竟该牺牲一致性还是可用性。
把这些零件放进真实的规模与压力中，它们会呈现出全新的形态。
因此，学习云计算不能只靠背诵概念或重复操作，而要在一个真实运行的系统中，把零散的知识重新组合并验证效果。

在本科计算机课程体系中，程序设计和数据结构帮助学生建立编写程序与组织数据的基本能力，计算机网络和操作系统进一步解释通信、进程、内存与 I/O 等基础机制。
进入真实在线系统后，这些知识不再彼此分离：一次请求既要经过网络传输，也要占用计算与存储资源，还可能与其他请求共同读写状态；规模扩大或节点发生故障后，系统还要处理排队、并发、数据一致性和服务恢复。
本书的定位，就在这些基础课程与复杂系统实践之间，帮助学生把已经学过的知识组织成分析、实现和改进在线系统的完整方法。

要把这个定位落到实处，需要一个看得见、摸得着的教学载体。
云计算本身太抽象，凭空讲，学生脑子里不容易有画面。
笔者在教学中尝试过不同的载体。
最早用过选课系统。
学期选课开始时，大量学生同时登录、查询和选课，本身就是一个真实的高并发场景；学生也熟悉选课的业务流程，不需要额外解释规则。
但选课系统的核心操作主要是数据的增删改查，状态冲突集中在"名额剩余量"这一处，很难自然引出分布式状态同步、跨机容错、弹性伸缩等更复杂的问题。
场景偏轻，学到的机制就有限。
后来又尝试过大模型服务系统。
推理服务的队列调度、批量处理、显存管理和横向扩容，几乎覆盖了云计算中资源调度与弹性伸缩的全部要点，场景足够真实，也足够前沿。
但它有两个门槛：一是依赖 GPU 硬件，不是所有学生的设备都能支撑；二是推理模型本身是一个黑盒，学生很难从底层理解请求在模型内部如何处理，也就难以完整追踪请求从进入到返回的全链路。
它适合拆解来学习局部机制，却不适合让学生从第一行代码开始亲手搭建完整系统。
经过反复比较，最终选择了在线游戏作为最终实验案例。

和选课系统相比，游戏的问题密度要高得多。
一场几十人的对战中，服务器要同时维持几十条长连接，每秒处理上百条操作指令，还要在多个玩家之间同步位置、状态和战斗结果。
高密度的资源并发、复杂的分布式状态同步、明显的弹性伸缩需求，这三类云计算的核心问题，在同一个游戏服务端上同时存在、彼此牵连。
学生在实现游戏的过程中，几乎每解决一个问题，就会自然遇到下一个问题：
连通了网络，就会遇到并发连接的排队；
并发跑起来了，就会碰到共享状态的冲突；
单机撑不住了，就得考虑拆分和同步；
节点多了以后，又要面对故障和治理。
问题一个推着一个往前走，学习也就有了连续的方向和动力。

和大模型服务系统相比，游戏的整条链路都是透明的。
从客户端发出操作，到服务器接收、校验、裁决、同步，再到结果返回，每一步都可以用代码亲手实现，每一处延迟也都可以通过工具测量和定位。
学生不需要依赖特殊硬件，也不需要面对推理模型这个黑盒，只要一台普通的笔记本电脑，就可以从第一行代码开始，把一个在线服务从最小原型逐步搭建到能够承载多人对战的完整系统。

游戏还有一个额外的优势：它足够直观，也足够有吸引力。
看得见摸得着，跑起来就能真人对战，操作是否流畅、延迟是否可接受，所有人都能立刻感受到。
更重要的是，为了让自己的操作响应再快十几毫秒，学生往往会主动去啃 epoll、一致性协议这些原本听起来枯燥的内容。
这种来自问题本身的学习动力，比任何外部考核都更稳定。

从最早的选课系统实验，到后来的大模型服务系统项目，再到如今的在线游戏案例，教学载体在近十年的课堂实践中一步步迭代演化。
每一次调整，都对应着对上一轮教学中暴露问题的修正；每换一个载体，都让问题链路更完整、学生动手的门槛更合适。
经过反复打磨的教学内容和方法，最终被系统地整理出来，就是面前这本书。

\section*{这本书有什么不一样}

本书既不是计算机网络、操作系统或分布式系统的专门教材，也不以某个云平台的产品和配置命令为主线。
它关注的是中间这一层：面对规模、状态、故障和资源约束时，怎样从具体现象识别系统问题，怎样选择合适的机制，并通过代码、实验和运行指标检验选择是否有效。
因此，本书既可以用于云计算和系统工程类课程，也可以作为课程设计、综合实践以及进一步研究云原生系统的基础。
具体来说，它有以下三个特点。

第一，以持续演进的在线系统案例引出问题，技术由真实约束推动。这本书从最小的双人对战原型起步，让系统在规模压力下一步步演进：先做出能跑的最小闭环，再处理单机并发；单机资源接近上限后，把服务拆到多台机器；机器和实例继续增加后，再引入容器与编排平台；最后进入基础设施的机制层，解释容器、网络和调度如何工作。每一章的技术都对应前一阶段暴露的具体问题。状态冲突发生后，才需要讨论一致性与共识；负载变化导致排队和延迟升高后，才需要讨论弹性伸缩。读者由此建立的不是技术名词之间的并列关系，而是问题、约束、机制、取舍和新问题之间的因果关系。

第二，把零散的知识组织成完整的系统判断。本书以在线游戏作为主要教学案例，同时使用在线协作、电商和大模型服务等系统呈现同类问题。案例可以不同，但背后的传输协议、并发控制、一致性、容器和调度机制具有共通性。设计消息协议时，计算机网络中的 TCP、UDP 有了具体作用；组织并发时，操作系统中的进程、线程和锁落到真实代码；拆分系统时，数据结构中的哈希方法成为分配对象与请求的重要工具。本书希望把原本分散在不同课程中的知识连接起来，使读者能够把在一个系统中形成的判断迁移到其他在线系统。

第三，实践材料与机制解释相互对应，内容进一步延伸到研究问题。配套案例提供可运行的代码，正文则通过必要的机制草图解释关键过程。全书从 TCP 与 UDP 的取舍、锁与 I/O 多路复用，走到一致性哈希、Raft 与 CAP，再走进容器、Kubernetes、HPA 与 Serverless 的启动与调度，最后讨论面向 AI Agent 的高并发 Sandbox。最后一章保留若干尚需继续探索的问题，为希望深入学习的读者提供进入云原生系统研究的起点。

云计算的水面之上，新技术来了又去；水面之下的基本原理，相对稳定。理解了原理，才能判断哪些变化只是表层，哪些真正触及结构。这也是本书不追热点、而从机制讲起的原因。

\section*{这本书如何学}

阅读本书不需要预先了解云计算或分布式系统的专门知识，但需要具备基本的编程能力：能够用一门语言写出百行规模的程序，理解数组、链表、哈希表等基本数据结构。如果接触过计算机网络或操作系统中的进程、内存、TCP 连接等概念，会更容易跟上。全书六章构成一条由问题驱动的主线，建议按顺序阅读：

\begin{itemize}
    \item 第一章 谋定全局：在线系统架构，拆解一个真实在线服务的全貌，建立评价架构的四个标尺，即性能、正确性、可用性、成本；
    \item 第二章 双雄集结：网络通信，写出第一个双人在线游戏，练成网络编程与消息协议的完整功底；
    \item 第三章 英雄集结：单机并发，讲如何组织执行、保护状态、复用资源，把一台机器的潜力榨干；
    \item 第四章 裂土封疆：分布式系统，分片、协同、容错三大主题一次讲透；
    \item 第五章 飞升入定：云原生部署，讲容器化交付、编排调度、弹性与自愈；
    \item 第六章 穷理尽微：云原生核心原理，深入容器、网络虚拟化与弹性调度的机制深处，并探索 Agent Sandbox 这一研究前沿。
\end{itemize}

对学习本书的读者，有三条操作上的建议。
其一，动手。每一章的代码都要亲手跑通，游戏必须真的对战起来，没跑通的章节等于没读。
其二，追问。每章开头先合上书，想一想"换我会怎么解决"，再与书中的方案比一比。两者的差距，通常就是收获最大的地方。
其三，善用 AI 而不依赖 AI。让 AI 回答具体的追问，而不是替你完成作业。能鉴别 AI 答案真伪的，恰恰是本书想建立的底层认知，两者互为补充。

对采用本书作为授课教材的教师：本书按一学期十六周设计，四个演进阶段各配一个可展示、可实战的课程项目（双雄对决、英雄集结、切分世界、飞升入定）；若课时不足，最后一章可压缩为一周导览，但前四章建议给学生留足动手时间。

本书与配套代码在 Gitee 开源、持续迭代（\texttt{https://gitee.com/hnu-cloudcomputing}），欢迎提交 issue 与 PR。授课视频、PPT 等配套资料，也可在上述仓库中找到。

使用这些材料时，建议把正文、配套代码、实验任务和章末思考题作为一个整体。
先运行给定系统并记录吞吐、延迟、资源占用或错误现象，再修改其中一个机制并比较结果，最后回到正文解释变化产生的原因。书稿与代码持续更新，实践时应尽量使用相互对应的版本，避免因接口或环境差异影响实验判断。

\subsection*{致谢}

本书由陈果、徐方林负责编写。
各章初稿撰写分工如下：胡文举（第1章）、庞海鑫（第2章）、谢先衍（第3章）、贺臻（第4章）、张道平（第5章）、徐方林（第6章）。
胡文举承担了课程实验内容的验证与复核，并参与指导了书中配图的视觉修订工作；终稿的核定由徐方林统一完成。

实验案例的设计与实现离不开历届学生的参与，陈俊杰、王煌等同学在游戏案例代码、实验文档和测试用例方面做了大量工作。
课程视频的制作同样有赖于学生团队的协助，刘筱芊等同学参与了授课视频的录制、剪辑与后期整理。

出版社的编辑老师们为本书的出版同样付出了大量心血，在此一并感谢。

受知识水平和时间精力所限，本书难免有错误和疏漏。
欢迎读者将发现的问题反馈至邮箱 \texttt{guochen@hnu.edu.cn}，我们将持续改进本书的质量。

\vspace{2em}

\begin{flushright}
陈果 2026年7月于长沙
\end{flushright}
